\documentclass[12pt]{article}
\usepackage[a4paper,margin=1in]{geometry}
\usepackage{newtxtext,newtxmath}      % Times New Roman look
\usepackage{setspace}\onehalfspacing  % 1.5 line spacing
\usepackage{array,booktabs,longtable}
\usepackage{enumitem}
\usepackage[table]{xcolor}
\usepackage{amsmath}
\usepackage{titlesec}
\usepackage[hidelinks]{hyperref}

\setlist{nosep,leftmargin=1.6em}
\titleformat{\section}{\large\bfseries}{\thesection.}{0.5em}{}
\renewcommand{\arraystretch}{1.25}
\newcommand{\g}[1]{\cellcolor{gray!45}#1}

\begin{document}

% ======================================================================
\begin{center}
{\large\bfseries Half-Yearly Progress Report for Jan--May 2026}\\[2pt]
{\bfseries Data sheet for MS scholars}
\end{center}

\begin{enumerate}[label=\arabic*.,leftmargin=2.2em]
\item \textbf{Name:} Ashutosh Patidar
\item \textbf{Registration No.:} DA24S006
\item \textbf{Department:} Data Science and AI
\item \textbf{Date of Joining:} 1st July 2024
\item \textbf{Specialization / Stream:} Data Science and AI
\item \textbf{Area of Research Work:} Interpretability (Transformers)
\item \textbf{Category of Admission (Regular/Part Time/External/Project\ldots):} Regular
\item \textbf{Guide:} Prof.\ Harish G.\ Ramaswamy
\item \textbf{Co-Guide(s) (if any):} ---
\item \textbf{Dates of GTC Meetings:}
\end{enumerate}

\begin{center}
\begin{tabular}{|p{3.4cm}|p{7.6cm}|p{2.4cm}|}
\hline
\textbf{Description} & \textbf{Event} & \textbf{Date}\\
\hline
1$^{\text{st}}$ GTC Meeting & Progress review / Research proposal meeting (Within 18 months from the date of registration) (Mandatory) & \\
\hline
2$^{\text{nd}}$ GTC Meeting & Research seminar (Within 27 months from the date of registration) (Mandatory) & \\
\hline
3$^{\text{rd}}$ GTC meeting & Thesis submission meeting (within 30 months from the date of registration) (Mandatory) & \\
\hline
Subsequent GTC meetings & Every 6 months till maximum period of the program or thesis submission whichever is earlier (Mandatory) & \\
\hline
\end{tabular}
\end{center}

\bigskip
\noindent\textbf{11. Details of course work}
\begin{center}
\begin{tabular}{|c|c|p{4.6cm}|c|c|c|}
\hline
\textbf{S.No} & \textbf{Course No.} & \textbf{Course Title} & \textbf{Sem.\ \& Year} & \textbf{Credit} & \textbf{Grade}\\
\hline
1 & DA5000 & Mathematical Foundations of Data Science & July--Nov 2024 & 12 & S\\
\hline
2 & DA5400 & Foundations of Machine Learning & July--Nov 2024 & 12 & B\\
\hline
3 & DA5401 & DA Lab & July--Nov 2024 & 6 & A\\
\hline
4 & DA6401 & Introduction to Deep Learning & Jan--May 2025 & 12 & A\\
\hline
5 & EE6180 & Advanced Topics in AI & Jan--May 2025 & 9 & A\\
\hline
6 & DA5402 & Machine Learning Operations Lab & Jan--May 2026 & 9 & S\\
\hline
\multicolumn{5}{|r|}{\textbf{OVERALL CGPA}} & 9.15\\
\hline
\end{tabular}
\end{center}

\bigskip
\noindent\textbf{12. Issues affecting research progress, if any:} None.

\vspace{1.4cm}
\noindent\begin{tabular}{p{5cm}p{5cm}p{5cm}}
\hrulefill & \hrulefill & \hrulefill\\
\textbf{Signature of Scholar} & \textbf{Signature of Co-Guide} & \textbf{Signature of Guide}
\end{tabular}

\newpage
% ======================================================================
\noindent\textit{Report should be in Times New Roman, 12 point, 1.5 spacing}

\section*{A. Short Summary (upto 3 pages only): Mandatory}

\noindent\textbf{Title of research work:} Interpretability of In-Context Learning: Analyzing
Causal Structure Acquisition in Minimal Transformer Models.

\medskip
\noindent\textbf{Problem Definition / Research Objectives (upto 250 words):}
The emergence of In-Context Learning (ICL) in Large Language Models presents a significant
interpretability challenge. Our research addresses the fundamental question: how do Transformer
architectures physically encode and utilise the causal structure of sequential data to make
predictions? Because analysing large models is computationally intractable, we adopt a
``minimal model'' approach---small, attention-centric Transformers trained on synthetic data
whose generative rules we control exactly. Our experiments confirm that such stripped-down
models reach near-optimal accuracy on unseen sequences, indicating they learn the underlying
transition/causal structure rather than memorising tokens; the core problem is to mathematically
decipher \emph{how} this is realised in the weights. The primary objective is to deconstruct the
optimisation dynamics of the Query--Key ($W_{QK}=W_Q^{\top}W_K$) and Output--Value
($W_{OV}=W_O W_V$) circuits---explaining why gradient descent converges to specific, structured
configurations that mirror the data's causal properties---and to compare these \emph{learned}
solutions against exactly \emph{hand-constructed} ``ideal'' circuits, using phase-wise training
and weight engineering to isolate each component. A second objective extends the analysis from a
single causal structure to the realistic regime of many competing in-context structures (random
Markov chains corrupted by noise), studying how a single model infers and applies the active rule
from context, and how noise density and architectural choices (attention parametrisation,
feed-forward layers) shape what gradient descent ultimately discovers---a step toward a rigorous,
mechanistic account of general-purpose In-Context Learning.

\medskip
\noindent\textbf{Brief review of literature (upto 250 words):}
Recent literature on Transformer interpretability has shifted from observing In-Context Learning
(ICL) performance to understanding the circuits that enable it. Following Brown et al.'s discovery
of ICL, Olsson et al.\ hypothesised that ``induction heads'' are the primary driving mechanism.
While theoretical studies established the expressivity of transformers---proving they can
implement algorithms such as gradient descent (Aky\"urek et al.; Von Oswald et al.) or Newton's
method (Fu et al.)---these works do not address whether gradient descent actually \emph{finds}
such solutions during training. Consequently, focus has shifted to the \emph{optimisation
dynamics} of ``tiny''/simplified transformers. Initial dynamical analysis was limited to single
attention layers solving simple tasks such as linear regression (Ahn et al.; Mahankali et al.).
Concurrent research now links these dynamics to the emergence of specific mechanistic circuits:
Bietti et al.\ and Reddy empirically demonstrate the sudden formation of induction heads during
training, while Edelman et al.\ and Aky\"urek et al.\ show that transformers trained on synthetic
languages undergo a sequential learning process, often mastering simpler unigram statistics before
developing the circuitry for bigrams or more complex causal structures. This body of work
collectively frames ICL not merely as statistical correlation, but as the learnable emergence of
algorithmic induction heads---motivating our mechanistic, construction-grounded study of how, and
under what conditions, these circuits form and what gradient descent converges to.

\medskip
\noindent\textbf{Research topic/gaps/tasks identified (upto 250 words):}
\emph{Topic:} mechanistic interpretability of the optimisation dynamics of small two-layer
``disentangled'' transformers performing in-context learning of \emph{noisy} Markov chains, and a
precise, weight-level comparison between gradient-descent solutions and exactly hand-constructed
induction circuits. \emph{Identified gaps:} (1) \textbf{Construction vs.\ learnability}---prior
work proves transformers \emph{can represent} induction/count circuits but rarely checks, weight
for weight, whether gradient descent \emph{converges} to the hand-constructible optimum, or to a
different solution related only by symmetry. (2) \textbf{Role of noise/structure}---existing
analyses use clean single-structure data; the effect of noise \emph{density} (a tunable signal
vs.\ noise mixture) on circuit formation, and the resulting \emph{phase transitions}, is
under-explored. (3) \textbf{Parametrisation and gauge}---attention is usually studied as separate
$W_Q,W_K,W_V,W_O$; the consequences of the combined $W_{QK}/W_{OV}$ parametrisation for
\emph{identifiability} (removing the within-head rotation gauge) and for \emph{trainability} are
not characterised. (4) \textbf{Feed-forward role}---the precise function of the position-wise FFN
in a near-attention-only induction model (calibration vs.\ computation) is not isolated.
\emph{Tasks:} (i) hand-engineer the exact induction circuit and the analytic optimal predictors
(context-count, ideal-Bayes); (ii) train GD models and compare cross-entropy/accuracy and weights
against the construction; (iii) explain the noise-driven phase transition; (iv) characterise the
combined parametrisation's identifiability/trainability trade-off; (v) isolate the FFN's role.

\medskip
\noindent\textbf{Summary of work done up to previous review (upto 250 words):}
In the previous period I studied mechanistic interpretability of ICL in ``tiny'' transformers
trained on synthetic \emph{periodic} data (e.g.\ \texttt{ABAB}, \texttt{ABCDE}). Using a
reductionist methodology, I progressively stripped the architecture---removing feed-forward
networks and LayerNorm, replacing learnable embeddings with one-hot (identity) embeddings, and
reducing the model to a single attention layer and head---to expose the minimal circuit that
solves cyclic next-token prediction. A key finding was that \textbf{positional embeddings are
strictly necessary}: without them self-attention is permutation-invariant over the context, so the
model cannot distinguish cyclic permutations (e.g.\ \texttt{ABC} vs.\ \texttt{BCA}) and fails to
generalise. To analyse how gradient descent reaches the solution, I built a custom
\emph{phase-wise} training framework that alternately freezes and updates the Query--Key
($W_{QK}$) and Output--Value ($W_{OV}$) circuits. This revealed a consistent two-stage dynamic:
$W_{OV}$ converges quickly to a task-optimal output map (sending each attended token to its correct
successor) \emph{largely independently} of $W_{QK}$, after which the $W_{QK}$ gradients are
``guided'' to place attention on the relevant previous position, driving the model to
$100\%$ accuracy; gradient tracking showed the simplified loss landscape is benign enough for GD to
reach the global optimum once the value map is established. Relating these findings to the
literature (Bietti et al.; Edelman et al.) confirmed alignment with theory for \emph{single}
causal structures and identified the central open problem---\emph{multiple competing causal
structures}---that motivated the current period's shift to a richer, noisy in-context Markov
setting.

\medskip
\noindent\textbf{Work done during the current review period (upto 250 words):}
This period operationalised the multi-structure goal as a controlled in-context Markov task: each
length-200 sequence carries its \emph{own} random transition table (each of $10$ signal states has
two equiprobable successors), and a tunable $\alpha$ mixes signal with noise---with probability
$\alpha$ the next token follows the in-context Markov rule, otherwise a uniform ``noise'' token
(a separate $10$-symbol vocabulary) is emitted. I built a two-layer ``disentangled'' transformer
(concatenating attention outputs onto a one-hot token/position stream) and, in parallel,
\emph{hand-engineered} the exact induction circuit (``$W^{*}$'') that implements the in-context
successor-frequency predictor, together with analytic baselines: the \emph{context-count}
predictor and the \emph{ideal-$\alpha$} Bayes floor. I trained GD models across
$\alpha=0.1\text{--}0.9$ under several heads (plain softmax; a learnable unembedding bias /
noise-floor; a freeze-readout variant; an ``oracle'' pipeline; a buffer construction) and compared
cross-entropy and accuracy against the construction and baselines. Key results: (i) GD essentially
\emph{matches} the context-count predictor, the residual gap to the Bayes floor being a
finite-context effect; (ii) a sharp \emph{noise-driven phase transition} governs whether GD learns
the hard noise-skipping circuit ($\alpha\!\approx\!0.8$ is hardest); (iii) collapsing attention
into the combined $W_{QK}/W_{OV}$ parametrisation \emph{removes the within-head gauge} (improving
identifiability) but converges markedly \emph{slower}, traced to the loss of the implicit
acceleration / low-rank bias of the over-parametrised factorisation; (iv) the position-wise FFN's
role is \emph{calibration}---a log-readout plus noise floor---that lifts the construction onto the
context-count optimum.

\medskip
\noindent\textbf{Future workplan (at least for the next 6 months), as Gantt chart or similar
chart:}
Over Jun--Nov 2026 I will: \textbf{T1} generalise from the two-successor toy to \emph{full} random
Markov chains (arbitrary, unequal transition probabilities) and quantify where a feed-forward
construction can beat plain context-counting; \textbf{T2} derive and verify when GD recovers the
hand construction \emph{up to} the gauge and temperature symmetries; \textbf{T3} design a
gauge-reduced yet trainable parametrisation (low-rank factored $W_{QK}/W_{OV}$) that keeps the
implicit low-rank bias while collapsing the gauge; \textbf{T4} extend the noise-driven
phase-transition and FFN-role analyses to multi-layer / multi-head settings; \textbf{T5}
consolidate results and prepare a manuscript.

\begin{center}
\begin{tabular}{|p{5.6cm}|c|c|c|c|c|c|}
\hline
\textbf{Task \textbackslash\ Month (2026)} & Jun & Jul & Aug & Sep & Oct & Nov\\
\hline
T1 Full Markov-chain generalisation & \g{} & \g{} & \g{} & & & \\
\hline
T2 GD-vs-construction guarantees & & \g{} & \g{} & \g{} & & \\
\hline
T3 Gauge-reduced low-rank parametrisation & & & \g{} & \g{} & \g{} & \\
\hline
T4 Multi-layer / multi-head extension & & & & \g{} & \g{} & \g{} \\
\hline
T5 Consolidation \& manuscript & & & & & \g{} & \g{} \\
\hline
\end{tabular}
\end{center}

\medskip
\noindent\textbf{Visible research output}
\begin{itemize}
\item Publications:
\item Conferences:
\item Workshops:
\end{itemize}

\newpage
% ======================================================================
\section*{B. Detailed Report (no page limit): Mandatory}

\section{Overview and the in-context Markov-noise task}
The previous review established, in tiny attention-only transformers on periodic data, that
positional embeddings are necessary and that the value circuit ($W_{OV}$) leads the query circuit
($W_{QK}$) during learning. The central open problem identified there was \emph{multiple competing
causal structures}. During this period I made this concrete with a single, controlled benchmark
that I use throughout (folders \texttt{ICL/1}--\texttt{ICL/8}).

\paragraph{Data-generating process.} The vocabulary has $V=20$ symbols: \emph{signal} tokens
$0\text{--}9$ and \emph{noise} tokens $10\text{--}19$. Each sequence (length $200$) draws its own
random transition table in which every signal state has exactly two equiprobable successors. A
threshold $\alpha$ controls the signal/noise mixture: with probability $\alpha$ the next token is a
successor of the current state (the state then advances), and with probability $1-\alpha$ a uniform
noise token is emitted (the state is unchanged). The model must \emph{infer the transition table
from context} and predict the next token; loss and accuracy are measured on the last $100$ tokens.
The table is resampled per sequence, so the task is genuine in-context learning of a fresh causal
structure each time, and $\alpha$ provides a continuous ``difficulty / noise-density'' knob.

\paragraph{Model.} A two-layer \emph{disentangled} transformer: each layer is a concat-attention
block whose output is appended to a residual stream initialised with one-hot token and positional
encodings; a linear unembedding (with optional bias) produces the logits. This keeps the
representation transparent (every coordinate is interpretable) while remaining expressive enough to
implement the induction circuit. Three reference predictors anchor every comparison: the
exactly hand-built circuit (\emph{$W^{*}$}, below); the \emph{context-count} predictor (the
empirical in-context successor frequencies, smoothed and $\alpha$-scaled, with the analytic noise
floor); and the \emph{ideal-$\alpha$} predictor (the Bayes floor that knows the true table).

\section{Weight engineering: the exact induction circuit ($W^{*}$) (\texttt{ICL/1})}
I first hand-constructed the weights of a transformer (\emph{PerfectInductionTransformer}) that
solves the task exactly, to obtain a ground-truth ``what a correct circuit looks like.'' Layer~1
implements a previous-token / positional matcher and value copy; layer~2 implements the induction
look-up: for the current state it attends to the positions that followed previous occurrences of
that state and averages their (one-hot) successors, yielding the in-context successor
\emph{frequency distribution}; the unembedding selects this block. A scalar temperature controls
how sharply the softmax reads out the frequencies (I observed that scaling the logits makes the
output more deterministic). This construction is the reference against which all trained models are
measured, and it doubles as a constructive proof that the two-layer disentangled architecture can
represent the optimal context predictor.

\section{Training with an unembedding bias / noise-floor (\texttt{ICL/2})}
The hand circuit outputs signal-token frequencies but must also place the correct mass
$(1-\alpha)/10$ on noise tokens. I therefore trained the softmax model with a \emph{learnable bias
at the unembedding}, giving gradient descent an explicit noise-floor degree of freedom. Using the
canonical data generator (10 signal / 10 noise symbols, two successors per state, length $200$), I
trained across $\alpha=0.1\text{--}0.9$ and stored per-threshold checkpoints and metrics. This
established the trained-model baseline and confirmed that a noise-floor parameter is needed for the
readout to represent the signal/noise split.

\section{Construction vs.\ trained-with-bias: where they differ (\texttt{ICL/3})}
I compared the hand construction against the GD-trained-with-bias model in cross-entropy. The
informative step was to \emph{split} the cross-entropy by target type---signal (``ruled'') tokens
vs.\ noise tokens---rather than reporting a single aggregate. This revealed that the apparent gap
was dominated by the \emph{noise-token} CE (which can explode if the floor is mis-set), and that
the bias indeed helps the trained model represent the floor. It also exposed a calibration error in
my early hand construction's probability distribution (the raw frequencies were not normalised into
a valid distribution before read-out), which I noted and later corrected. This step set the
methodology of \emph{always decomposing error by token type and by context} that I use throughout.

\section{The oracle training pipeline (\texttt{ICL/4})}
I consolidated the experiment into a reproducible pipeline: a single \texttt{data\_generation.py}
(canonical seeds, train/test splits, per-$\alpha$ files), a \texttt{train.py} training the
disentangled transformer with the loss and evaluation restricted to the last $100$ tokens, and a
\texttt{dataset/} and \texttt{models/} layout reused by all later folders. ``Oracle'' refers to
training and measuring against the Bayes-optimal target (ideal-$\alpha$) so that every trained
model's cross-entropy can be read as a gap to the information-theoretic floor. This pipeline (with
the attention width matched to the one-hot construction) is the backbone for folders
\texttt{ICL/6}--\texttt{ICL/8}.

\section{Buffer construction: making the $\alpha$-predictor implementable (\texttt{ICL/5})}
To make the $\alpha$-scaled count predictor exactly realisable inside the architecture, I
introduced a \emph{buffer} construction---auxiliary residual capacity that holds the intermediate
counts/normalisation so the readout can produce the $\alpha$-implementable distribution. The
analysis verified the ``$\alpha$-implementable construction'': with the buffer, the network can
emit the context-count distribution (signal frequencies scaled by $\alpha$ plus the noise floor)
rather than the distorted output a bare softmax would give. This is the constructive precursor to
the later finding that a feed-forward block can play the same calibrating role.

\section{Understanding what gradient descent figures out (\texttt{ICL/6})}
This folder is the analytical core. Across $\alpha=0.1\text{--}0.9$ and on both splits I compared,
in one framework, the GD softmax model, the $W^{*}$ construction (at its best temperature), the
context-count predictor, and the ideal-$\alpha$ floor.

\paragraph{GD tracks context-counting.} The trained softmax model's cross-entropy closely follows
the context-count predictor; the residual gap to ideal-$\alpha$ is a finite-context effect (the
empirical successor frequencies are noisy estimates of the true two-successor distribution).

\paragraph{A noise-driven phase transition.} The genuinely hard sub-task is predicting
\emph{immediately after a noise token}, where the model must skip the noise and recover the state
from the last signal token. Measuring the post-noise wrong-successor rate showed a sharp transition:
it is essentially zero for $\alpha\le 0.7$ but jumps to $\sim 50\%$ for $\alpha\ge 0.8$. The reason
is that the incentive to learn the costly noise-skipping circuit scales with the noise density
$1-\alpha$; once noise is rare, gradient descent under a fixed budget does not reliably form the
circuit, so $\alpha\approx 0.8$ is the worst case (past the transition yet still $20\%$ noise). I
corroborated this with a freeze-readout variant (the readout pinned to $W^{*}$ with only a learnable
noise-floor bias) and with per-position diagnostics that display, for a single query, every
predictor's distribution alongside the in-context evidence (how many times each candidate successor
was actually seen following the current state).

\section{Combined $W_{QK}/W_{OV}$ parametrisation and gauge (\texttt{ICL/7})}
Attention uses $W_Q,W_K$ only through the product $W_{QK}=W_Q^{\top}W_K$ and $W_V,W_O$ only through
$W_{OV}=W_OW_V$. I therefore re-parametrised both layers directly in terms of $W_{QK}$ and $W_{OV}$
(\texttt{train\_combined.py}, plain softmax, \emph{no} feed-forward), which removes the within-head
$\mathrm{GL}(d)$ rotation gauge and shrinks the solution orbit---the motivation being
\emph{identifiability}, so that trained solutions can be meaningfully compared. I then
\emph{collapsed} the hand construction into the same parametrisation (verified to reproduce its
logits exactly, $\max|\Delta|=0$) and asked: does GD converge to the construction? Findings:
(i) functionally, the trained model and the collapsed construction agree only \emph{up to a
temperature} and the residual softmax / readout gauges---so they implement the same computation in a
different gauge rather than identical weights; (ii) crucially, the combined parametrisation, while
more identifiable, \emph{converges substantially slower} than the separate $W_Q,W_K,W_V,W_O$ form
(it needed several times more epochs to reach the same accuracy). I traced this to the
\emph{implicit acceleration} of the over-parametrised factorisation: gradient descent on a product
$M=AB$ preconditions the update on $M$ by the factor norms, descending much faster than gradient
descent on a single matrix $M$. This is an identifiability-vs-trainability trade-off central to the
project.

\section{The role of the feed-forward block (\texttt{ICL/8})}
Finally I studied adding position-wise feed-forward (FFN) blocks to the disentangled model: a
minimal single-FFN model and a ``two-cell'' model with an FFN after \emph{both} attention layers,
each in the separate and the combined parametrisation. Results: (i) a trained FFN markedly improves
calibration and lands between context-count and the ideal floor; (ii) the FFN's essential role is a
\emph{log-readout plus noise floor}---a plain softmax of the attention's count vector
\emph{exponentiates} it and is distorted, whereas an FFN that computes $\log(\text{count}+\text{floor})$
lets the softmax recover the normalised counts, i.e.\ exactly the context-count optimum. I verified
this by hand-constructing the FFN as $\log(L+b)$ with separate signal/noise floors, which reproduces
context-count to within $0.01\text{--}0.03$ nats and clearly beats the fixed-temperature softmax of
the same circuit. (iii) The achievable floor for any context-only model is the
\emph{context-Bayesian} predictor, which sits just below context-count and just above ideal-$\alpha$;
the FFN cannot pass ideal-$\alpha$ because the true table is not fully revealed by a finite context.
(iv) Ablation/linearisation showed that, in the \emph{trained} two-cell model, both FFNs do
load-bearing nonlinear work (removing or linearising either one sharply raises CE), with a
$\alpha$-dependent division of labour (the first FFN matters most in the high-noise regime, the last
FFN dominates the final calibration). Together with \texttt{ICL/6}--\texttt{ICL/7}, this gives a
coherent mechanistic account of \emph{what} the trained network computes, \emph{why} it sometimes
fails to (noise-driven phase transition), and \emph{how close} it can come to the optimum (the
log-readout calibration and the context-Bayes ceiling).

\section{Summary of contributions this period}
\begin{itemize}
\item A controlled in-context \emph{noisy-Markov} benchmark with a tunable noise knob $\alpha$, and
a transparent two-layer disentangled transformer for studying it.
\item An exact hand-engineered induction circuit ($W^{*}$) plus analytic optimal predictors
(context-count, ideal-$\alpha$, context-Bayes) used as a measuring stick for every trained model.
\item A mechanistic explanation of a \emph{noise-driven phase transition} in circuit formation
(why $\alpha\approx 0.8$ is hardest).
\item A characterisation of the combined $W_{QK}/W_{OV}$ parametrisation: it removes the within-head
gauge (better identifiability) at the cost of slower convergence (loss of the factorisation's
implicit acceleration).
\item An isolation of the feed-forward block's role as a log-readout + noise-floor \emph{calibrator}
that realises the context-count optimum, together with the context-Bayes achievability ceiling.
\end{itemize}

\section*{References}
\begin{enumerate}[label={[\arabic*]},leftmargin=2.4em]
\item K.~Ahn, X.~Cheng, H.~Daneshmand, S.~Sra. Transformers learn to implement preconditioned
gradient descent for in-context learning. \emph{NeurIPS}, 2023.
\item E.~Aky\"urek, D.~Schuurmans, J.~Andreas, T.~Ma, D.~Zhou. What learning algorithm is
in-context learning? Investigations with linear models. \emph{ICLR}, 2023.
\item E.~Aky\"urek et al. In-context language learning: architectures and algorithms.
\emph{ICML}, 2024.
\item A.~Bietti, V.~Cabannes, D.~Bouchacourt, H.~J\'egou, L.~Bottou. Birth of a Transformer: a
memory viewpoint. \emph{NeurIPS}, 2023.
\item T.~B.~Brown et al. Language models are few-shot learners. \emph{NeurIPS}, 2020.
\item B.~L.~Edelman, E.~Edelman, S.~Goel, E.~Malach, N.~Tsilivis. The evolution of statistical
induction heads: in-context learning Markov chains. \emph{NeurIPS}, 2024.
\item N.~Elhage et al. A mathematical framework for transformer circuits.
\emph{Transformer Circuits Thread}, Anthropic, 2021.
\item D.~Fu, T.~Chen, R.~Jia, V.~Sharan. Transformers learn higher-order optimization methods for
in-context learning. 2023.
\item A.~Mahankali, T.~B.~Hashimoto, T.~Ma. One step of gradient descent is provably the optimal
in-context learner with one layer of linear self-attention. \emph{ICLR}, 2024.
\item E.~Nichani, A.~Damian, J.~D.~Lee. How transformers learn causal structure with gradient
descent. \emph{ICML}, 2024.
\item C.~Olsson et al. In-context learning and induction heads.
\emph{Transformer Circuits Thread}, Anthropic, 2022.
\item G.~Reddy. The mechanistic basis of data dependence and abrupt learning in an in-context
classification task. \emph{ICLR}, 2024.
\item J.~von Oswald et al. Transformers learn in-context by gradient descent. \emph{ICML}, 2023.
\end{enumerate}

\end{document}
